\section*{Asesmen}
\addcontentsline{toc}{section}{Asesmen}

Asesmen Bab IV mencakup Kuis Konsep (3\%) dan UTS Praktik (14\%).
Fokusnya adalah ketepatan tabel pemetaan Tipe 3--0, klasifikasi lima
grammar, dan hubungan teori dengan fase kompilasi.

\begin{table}[H]
\centering
\caption{Instrumen asesmen Bab IV}
\label{tab:ch04-asesmen}
\begin{tabularx}{\textwidth}{@{}lXl@{}}
\toprule
\textbf{Instrumen} & \textbf{Indikator} & \textbf{Bobot} \\
\midrule
Kuis Konsep & Definisi formal, hierarki Chomsky, dan jenis automata &
3\% (bagian) \\
UTS Praktik & Tabel pemetaan Tipe 3--0 dan klasifikasi lima grammar
beserta alasannya & 14\% (bagian) \\
Lokakarya & Tabel, klasifikasi, dan hasil uji lima grammar &
prasyarat \\
\bottomrule
\end{tabularx}
\end{table}

Kriteria kelulusan: mahasiswa dapat mengisi tabel pemetaan Tipe 3--0
ke automata, mengklasifikasikan lima grammar dengan alasan yang tepat,
dan menjelaskan mengapa setiap fase kompilator memerlukan tingkat
formalisme yang berbeda.
